% Context from chapters/optim-nonsmooth.tex; for mathematical reference, not compilation.
\section{Proximal Algorithm}

For a nonsmooth objective $f$, an explicit gradient step may be undefined. An implicit step remains available even when $f$ is nonsmooth.


                                                
\subsection{Proximal Map}

For a step size $\tau>0$, define the proximal map by
\eql{\label{eq-defn-proximal}
	\foralls x, \quad
	\Prox_{\tau f}(x) \eqdef \uargmin{x'} \frac{1}{2}\norm{x-x'}^2+\tau f(x').
}
It balances a decrease in $f$ against a quadratic penalty for moving away from $x$.
 
Since $f\in\Ga_0(\Hh)$, the objective $\frac12\norm{x-\cdot}^2+\tau f$ is strongly convex and coercive, so $\Prox_{\tau f}$ is well defined and single-valued.

For an indicator $f=\iota_\Cc$, the proximal map is the projection $\Prox_{\iota_\Cc}=\Proj_{\Cc}$. It therefore generalizes projection to convex functions. For each fixed input, the proximal point is also the projection onto a suitable sublevel set of $f$. Like an orthogonal projection, the proximal map is nonexpansive.

\begin{prop}
	One has $\norm{ \prox_{f}(x)-\prox_{f}(y) } \leq \norm{x-y}$.
\end{prop}



\begin{figure}[tbp]
\centering
\BookDrawing{tikz/optim-nonsmooth-prox-projection}
\caption{Proximal map and projection map.}
\label{fig-prox-proj}
\end{figure}


   
\paragraph{Examples}

Several common proximal maps have closed forms.

\begin{prop}
One has
\eql{\label{eq-prox-l2-l1}
	\Prox_{\frac{\tau}{2} \norm{\cdot}^2}(x) =  \frac{x}{1+\tau}, 
	\qandq
	\Prox_{\tau\norm{\cdot}_1}(x)=\Ss_\tau^1(x),
}
where the soft-thresholding is defined as
\eq{
	\Ss_\tau^1(x) \eqdef ( S_\tau(x_i) )_{i=1}^p \qwhereq
	S_\tau(r)\eqdef\sign(r)(|r|-\tau)_+,
}
(see also~\eqref{eq-soft-thresh}). For $A \in \RR^{P \times N}$, one has
\eql{\label{eq-prox-quad}
	\Prox_{\frac{\tau}{2}\norm{A \cdot-y}^2}(x) = (\Id_N + \tau A^* A)^{-1} ( x + \tau A^* y ).
	                                                                 
}
\end{prop}
\begin{proof}
	The proximal map of $\norm{\cdot}_1$ was derived informally in Proposition~\ref{prop-equiv-sparse-thresh}.
	 
	For a rigorous proof, use separability to reduce to a scalar input $r$. The optimality condition for its proximal point $z$ is $0 \in z-r + \tau \partial |\cdot|(z)$.
	The candidate $z=0$ is optimal precisely when $0 \in -r + [-\tau,\tau]$, or equivalently $r \in [-\tau,\tau]$.
	 
	For $z \neq 0$, optimality requires $z = r - \tau \sign(z)$. When $r \notin [-\tau,\tau]$,
	the candidate $z = r - \tau \sign(r)$ has the same sign as $r$, so it is the unique minimizer.
	 
	For the quadratic case
	\eq{
		z = \Prox_{\frac{\tau}{2}\norm{A \cdot-y}^2}(x) 
		\quad\Leftrightarrow\quad 
		z-x + \tau A^*( A z - y ) = 0
		\quad\Leftrightarrow\quad
		(\Id_N + \tau A^* A) z = x + \tau A^* y.
	}
\end{proof}


For some nonconvex functions, the proximal minimization also has solutions, but they may be nonunique. For example,
$\Prox_{\tau \norm{\cdot}_0}$ is hard thresholding at $\sqrt{2\tau}$; see Proposition~\ref{prop-equiv-sparse-thresh}. At $|x_i|=\sqrt{2\tau}$, both $0$ and $x_i$ are minimizers.


                                          
\subsection{Basic Properties}

The following identities simplify proximal computations.

\begin{prop}
One has
\eq{
	\Prox_{f+\dotp{y}{\cdot}}=\Prox_f(\cdot-y), \quad
	\Prox